\section{Reward Shaping y estabilización del valor}

\subsection{Técnicas de ajuste fino para la estabilidad de la función de valor}

En las actualizaciones off-policy, el drift en la predicción de la función Q afecta gravemente al actor. Los stabilizers comunes en la práctica incluyen:

\begin{itemize}
    \item \textbf{Target smoothing}: en el TD target se usa twin Q average + Polyak average smoothing coefficient $\tau$.
    \item \textbf{Reward clipping}: limita el reward a $[-1,1]$ para evitar un gradient explosivo.
    \item \textbf{Value normalization}: en batch_level se realiza una normalización poblacional de los valores Q, para reducir el scale drift.
\end{itemize}

\subsection{Principios de diseño del Shaping}

El reward shaping debe seguir el potential-based shaping theorem, para garantizar que el shaping term no cambie la política óptima:

$$F(s, a, s') = \gamma \Phi(s') - \Phi(s)$$

donde $\Phi$ es la potential function. Un ejemplo práctico de $\Phi(s)$ es distance-to-goal, que mejora la guiding signal del agent en las early stages.

\begin{warningbox}{Consecuencias de un shaping incorrecto}
Si el shaping recompensa directamente «acercarse al objetivo» en lugar de «completar la tarea», el agent puede quedarse antes de tiempo en un estado con un reward potencial alto, lo cual produce reward hacking. Por eso el potential-based shaping debe apegarse estrictamente a la MDP structure.
\end{warningbox}

\subsection{Resumen de la sección}

El reward shaping y la value stabilization son el núcleo para mantener estable el long-term training. Usar en conjunto target smoothing, reward clipping y potential-based shaping permite evitar eficazmente la value explosion y el reward hacking.

